\section{Why sealed loading was chosen here}
\label{app:vented}

Sealed loading is an architectural choice for this application, not the
outcome of a general topology contest. This appendix states the reasoning
and its limits.

A vented alignment is a fourth-order system \cite{small1973vented}: the box
compliance and the port's acoustic mass and resistance add a degree of
freedom to the sealed operator \cref{eq:ode}, to which it reduces exactly
as the port mass and resistance grow without bound. Compared with a sealed
box of the same volume, it buys roughly an octave of extension and higher
passband sensitivity---a decisive advantage when amplifier power, not
displacement or room gain, is the binding constraint.

Four properties make it the less convenient choice for the architecture of
\cref{sec:intro-arch}.

\textbf{Below-tuning excursion.} Below the tuning frequency the enclosure
supplies little restoring force, so cone excursion is limited by the drive
signal rather than by the alignment. Standard practice adds a protective
high-pass filter below tuning, which is entirely workable, but it removes
output in exactly the region where the pressure-zone branch of
\cref{eq:Vdem} would otherwise be exploited. The sealed alternative bounds
displacement at every frequency down to DC without an additional protective
element, which makes the excursion gates of \cref{sec:risk-policy}
sufficient on their own.

\textbf{Steeper low-frequency slope.} The vented low-frequency asymptote is
steeper than the sealed $12$\,dB/oct, so the approximate complementarity
with the room's pressure-zone branch \cref{eq:Fsrule} does not carry over
directly, and correction toward the bottom of the band starts from a larger
deficit.

\textbf{Port behaviour.} Near tuning most of the output passes through the
port, and the port air velocity at a given target scales as
$v_{\mathrm{port}}\approx2\pi r_{\mathrm{listen}}p_t/
(\rho_0\omega_b S_{\mathrm{port}})$. At the levels and frequencies of this
scenario a small port would run well above the velocity at which
compression, turbulence noise, and nonlinear behaviour become significant
\cite{salvatti2002,vanderkooy1998}. This is a design constraint on port
geometry, not a defect of vented loading: large flared ports, several
ports, or a passive radiator address it, at the cost of enclosure volume,
baffle area, or an additional moving part. A passive radiator removes the
port mechanism entirely while retaining the fourth-order alignment and the
below-tuning excursion behaviour.

\textbf{System integration.} The manifold of \cref{sec:intro-arch} already
uses several active drivers for displacement and first-order
reaction-force cancellation under matched opposed-pair operation, each
with its own amplifier channel. In that context the extra tuning element,
the protective filter, and the port geometry constraint add system
complexity without addressing the binding constraint, which
\cref{app:example} shows is enclosure volume and excursion rather than
amplifier power.

None of this establishes that sealed loading is generally superior. Under a
different constraint set---limited amplifier power, limited enclosure
volume with a higher extension requirement, or a room without usable
pressure-zone gain---a vented or passive-radiator alignment is the
better-matched choice, and the gates of \cref{sec:risk-policy} would have
to be re-derived for its fourth-order model.
